% Part: first-order-logic
% Chapter: models-theories
% Section: expressing-props-of-structures

\documentclass[../../../include/open-logic-section]{subfiles}

\begin{document}

\olfileid[hi]{fol}{mat}{exs}

\olsection{\printtoken{P}{structure}: उनके गुणधर्म व्यक्त करना}

\begin{explain}
फलनों और संबंधों पर शर्तें व्यक्त करना, अथवा अधिक सामान्यतः यह व्यक्त करना कि
किसी संरचना के फलन और संबंध उन शर्तों को संतुष्ट करते हैं, प्रायः उपयोगी और
महत्त्वपूर्ण होता है। उदाहरणतः किसी भाषा में कुछ !!{structure}s ऐसी होती हैं
जिनमें !!{predicate}s का वही ``अर्थ'' होता है जो हम उनसे चाहते हैं; हम उन्हें
अन्य संरचनाओं से अलग पहचानना चाहेंगे। निस्संदेह हम यह बताने के लिए पूर्णतः
स्वतंत्र हैं कि कौन-सी !!{structure}s हमारे ``अभीष्ट'' हैं। जैसे, हम कह सकते
हैं कि !!{predicate}~$\le$ की व्याख्या कोई क्रम होनी चाहिए; अथवा हमारी रुचि
केवल~$\Lang L$ की उन व्याख्याओं में है जिनका प्रांत समुच्चयों से बना है और
जिनमें~$\Obj \in$ की व्याख्या ``का !!a{element} है'' संबंध द्वारा होती है।
पर क्या हम यह काम भाषा के !!{sentence}s से कर सकते हैं? दूसरे शब्दों में,
!!a{structure}~$\Struct M$ पर कौन-सी शर्तें हम~$\Struct M$ की भाषा के किसी
!!a{sentence} (या संभवतः !!{sentence}s के किसी समुच्चय) द्वारा व्यक्त कर सकते
हैं? कुछ शर्तें ऐसी हैं जिन्हें हम व्यक्त नहीं कर सकेंगे। उदाहरणतः~$\Lang L_A$
का कोई ऐसा वाक्य नहीं जो किसी !!{structure}~$\Struct M$ में केवल
तभी सत्य हो जब $\Domain M = \Nat$। हम ``प्रांत में केवल प्राकृतिक संख्याएँ
हैं'' व्यक्त नहीं कर सकते। किंतु कुछ !!{structure}s शायद ऐसे ``संरचनात्मक
गुणधर्म'' रखती हैं जिन्हें व्यक्त किया जा सकता है। उन्हें रखने वाली
!!{structure}s कौन-सी हैं, और कौन-से गुणधर्म हम !!{sentence}s द्वारा व्यक्त कर
सकते हैं? अथवा दूसरे रूप में पूछें: कौन-सी !!{structure}s किसी !!a{sentence}
(या !!{sentence}s के समुच्चय) को सत्य बनाती हैं, और इस प्रकार संरचनाओं के किन
संग्रहों का वर्णन किया जा सकता है?
\end{explain}

\begin{defn}[समुच्चय का प्रतिरूप]
मान लें कि~$\Gamma$, भाषा~$\Lang L$ के !!{sentence}s का समुच्चय है। हम कहते
हैं कि !!a{structure}~$\Struct M$,~$\Gamma$ \emph{का प्रतिरूप है} यदि
$\Sat{M}{!A}$ सभी $!A \in \Gamma$ के लिए।
\end{defn}

\begin{ex}
वाक्य~$\lforall[x][x \le x]$,~$\Struct M$ में तभी और केवल तभी सत्य
है जब $\Assign{\le}{M}$ कोई स्वतुल्य संबंध है। वाक्य
$\lforall[x][\lforall[y][((x \le y \land y \le x) \lif x = y)]]$,
~$\Struct M$ में तभी और केवल तभी सत्य है जब $\Assign{\le}{M}$ प्रति-सममित
है। वाक्य~$\lforall[x][\lforall[y][\lforall[z][((x \le y \land y \le z)
      \lif x \le z)]]]$,~$\Struct M$ में तभी और केवल तभी सत्य है जब
$\Assign{\le}{M}$ संक्रामक है। अतः
\begin{align*}
\{\quad &\lforall[x][x \le x], \\
   & \lforall[x][\lforall[y][((x \le y \land y \le
    x) \lif x = y)]], \\
   &\lforall[x][\lforall[y][\lforall[z][((x \le y
      \land y \le z) \lif x \le z)]]] \quad \}
\end{align*}
के प्रतिरूप ठीक वे संरचनाएँ हैं जिनमें~$\Assign{\le}{M}$ स्वतुल्य,
प्रति-सममित और संक्रामक, अर्थात् आंशिक क्रम है। इसलिए हम इन्हें
\emph{आंशिक क्रमों के प्रथम-क्रम सिद्धांत} के अभिगृहीत मान सकते हैं।
\end{ex}

\end{document}
